% BEGIN VOL07-EXPANSION VII01
\section{Supplementary worked examples}
\label{sec:vii01-pedagogy}

\begin{example}[Two-chart atlas on the circle]\label{ex:vii01-ped-01}
Let \(U_N=S^1\setminus\{(0,1)\}\) and \(U_S=S^1\setminus\{(0,-1)\}\).  Stereographic projection identifies each set with \(\mathbb R\).  The overlap is \(S^1\) with two points deleted, so the two coordinates give an explicit topological atlas.  This separates the local Euclidean statement from the global fact that no single chart can cover the compact circle.
\end{example}

\begin{example}[A quotient chart on the torus]\label{ex:vii01-ped-02}
Write \(T^2=\mathbb R^2/\mathbb Z^2\).  If \(p\in\mathbb R^2\), choose a disc \(D\) of radius less than \(1/2\).  Distinct integer translates of \(D\) are disjoint, so the quotient map restricts to a homeomorphism \(D\to q(D)\).  Thus the quotient description produces the same local Euclidean structure as the product \(S^1\times S^1\).
\end{example}

\begin{example}[A local invariant detects a nonmanifold point]\label{ex:vii01-ped-03}
The wedge \(S^1\vee S^1\) is locally a line away from the wedge point.  Removing the wedge point from a sufficiently small neighborhood leaves four components, while deleting a point from an interval leaves two.  The number of local branches is therefore enough to rule out a \(1\)-manifold chart at the wedge point.
\end{example}

\section{Graded supplementary exercises}
The following exercises add a deliberate progression from concrete calculation to structural proof, hypothesis testing, application, and synthesis.

\subsection*{Standard computations and constructions}

\begin{exercise}[Circle charts]\label{exr:vii01-ped-01}
Write explicit stereographic coordinates on \(S^1\setminus\{(0,1)\}\).
\end{exercise}
\begin{hint}
Project from the missing point to the horizontal axis.
\end{hint}
\begin{solution}
If \((x,y)\in S^1\) with \(y\neq1\), one convenient coordinate is \(t=x/(1-y)\); the inverse is rational in \(t\) and gives a homeomorphism with \(\mathbb R\).
\end{solution}

\begin{exercise}[Product dimension]\label{exr:vii01-ped-02}
Construct a chart on \(S^1\times\mathbb R\) and identify its dimension.
\end{exercise}
\begin{hint}
Take the product of a circle chart with the identity chart on \(\mathbb R\).
\end{hint}
\begin{solution}
The product chart lands in an open subset of \(\mathbb R^2\), so the cylinder is a \(2\)-manifold.
\end{solution}

\begin{exercise}[Quotient neighborhood]\label{exr:vii01-ped-03}
For \(\mathbb R/\mathbb Z\), find an interval around \(0\) on which the quotient map is injective.
\end{exercise}
\begin{hint}
Use an interval of length strictly less than \(1\).
\end{hint}
\begin{solution}
For example \((-1/3,1/3)\) meets no nontrivial integer translate of itself, so it maps homeomorphically to its image.
\end{solution}

\begin{exercise}[Punctured disc]\label{exr:vii01-ped-04}
Determine the number of connected components of \(B_\varepsilon(0)\setminus\{0\}\subset\mathbb R^2\).
\end{exercise}
\begin{hint}
Use polar coordinates.
\end{hint}
\begin{solution}
It is path connected: any two points can be joined radially to a common small circle and then along that circle.
\end{solution}

\begin{exercise}[Too many components]\label{exr:vii01-ped-05}
Explain why an uncountable disjoint union of copies of \(S^1\) is not second countable.
\end{exercise}
\begin{hint}
Each component is open, and a countable basis must meet every nonempty open component.
\end{hint}
\begin{solution}
Distinct open components require distinct basis elements contained in them, so a countable basis cannot serve uncountably many components.
\end{solution}

\subsection*{Proofs}

\begin{exercise}[Open submanifolds]\label{exr:vii01-ped-06}
Prove that every open subset of an \(n\)-manifold is an \(n\)-manifold.
\end{exercise}
\begin{hint}
Restrict charts and a countable basis.
\end{hint}
\begin{solution}
Hausdorffness and second countability pass to open subspaces, and chart restrictions give local Euclidean neighborhoods of the same dimension.
\end{solution}

\begin{exercise}[Countable atlas]\label{exr:vii01-ped-07}
Prove that a second-countable locally Euclidean space has a countable atlas.
\end{exercise}
\begin{hint}
Refine chart domains by elements of one countable basis.
\end{hint}
\begin{solution}
For each point choose a basis element contained in a chart domain.  Only countably many basis elements occur, and the corresponding restricted charts cover the space.
\end{solution}

\begin{exercise}[Product theorem]\label{exr:vii01-ped-08}
Prove that the product of an \(m\)-manifold and an \(n\)-manifold is an \((m+n)\)-manifold.
\end{exercise}
\begin{hint}
Take products of local charts and countable bases.
\end{hint}
\begin{solution}
Product charts identify neighborhoods with opens in \(\mathbb R^{m+n}\); Hausdorffness and second countability are preserved by finite products.
\end{solution}

\begin{exercise}[Circle as a quotient]\label{exr:vii01-ped-09}
Prove that \(\mathbb R/\mathbb Z\) is homeomorphic to \(S^1\).
\end{exercise}
\begin{hint}
Use \(t\mapsto e^{2\pi i t}\) and the universal property of the quotient.
\end{hint}
\begin{solution}
The exponential map is constant exactly on integer cosets and descends to a continuous bijection \(\mathbb R/\mathbb Z\to S^1\); compactness of the quotient fundamental interval and Hausdorffness of \(S^1\) give a homeomorphism.
\end{solution}

\subsection*{Counterexamples and hypothesis tests}

\begin{exercise}[Local Euclidean implies Hausdorff]\label{exr:vii01-ped-10}
Test the claim that every locally Euclidean space is Hausdorff.
\end{exercise}
\begin{hint}
Use the plane with two origins.
\end{hint}
\begin{solution}
The claim is false: that space is locally Euclidean but its two origins cannot be separated by disjoint neighborhoods.
\end{solution}

\begin{exercise}[Quotients are automatically manifolds]\label{exr:vii01-ped-11}
Test the claim that every quotient of a manifold by a group action is a manifold.
\end{exercise}
\begin{hint}
Consider a nonfree action with a fixed point.
\end{hint}
\begin{solution}
False.  For example reflection of \(\mathbb R\) about \(0\) has a fixed point; quotient behavior near the fixed orbit differs from a free covering-type quotient.
\end{solution}

\begin{exercise}[A wedge is a one-manifold]\label{exr:vii01-ped-12}
Test whether \(S^1\vee S^1\) is a topological \(1\)-manifold.
\end{exercise}
\begin{hint}
Delete the wedge point from a small neighborhood.
\end{hint}
\begin{solution}
It is not: four local branches remain, whereas deleting the center of an interval leaves exactly two.
\end{solution}

\subsection*{Applications and investigations}

\begin{exercise}[Torus from a lattice]\label{exr:vii01-ped-13}
Use local quotient charts to explain why \(\mathbb R^2/\mathbb Z^2\) is a surface.
\end{exercise}
\begin{hint}
Choose discs smaller than half the lattice spacing.
\end{hint}
\begin{solution}
Such a disc has disjoint nontrivial translates, so the quotient map is locally a homeomorphism; Hausdorffness and second countability follow from the standard lattice quotient.
\end{solution}

\begin{exercise}[Covering-space diagnostic]\label{exr:vii01-ped-14}
Explain why a covering space of an \(n\)-manifold is locally an \(n\)-manifold.
\end{exercise}
\begin{hint}
Use an evenly covered neighborhood.
\end{hint}
\begin{solution}
Each sheet over an evenly covered chart is homeomorphic to that chart domain, hence locally Euclidean of the same dimension; the usual covering hypotheses supply the global separation/countability conditions in standard cases.
\end{solution}

\subsection*{Challenge problems}

\begin{exercise}[Local branch invariant]\label{exr:vii01-ped-15}
Construct a connected Hausdorff space that is locally a \(1\)-manifold except at one point with three local branches, and prove the exceptional point is not manifold-like.
\end{exercise}
\begin{hint}
Glue three half-lines at their endpoints.
\end{hint}
\begin{solution}
The resulting triod is connected and Hausdorff.  Removing the central point from a small neighborhood leaves three components, impossible for a punctured interval.
\end{solution}

\begin{exercise}[Compactness versus one chart]\label{exr:vii01-ped-16}
Show that a nonempty compact \(n\)-manifold cannot be covered by a single chart whose image is all of \(\mathbb R^n\).
\end{exercise}
\begin{hint}
A homeomorphism preserves compactness.
\end{hint}
\begin{solution}
If such a chart existed, the compact manifold would be homeomorphic to noncompact \(\mathbb R^n\), a contradiction.
\end{solution}

% END VOL07-EXPANSION VII01
